\documentclass[sigconf,nonacm]{acmart}
\usepackage{booktabs}
\begin{document}

\title{Inference, Not Just Notation:\\ Notes on Sampling-Based
  Key Extraction for Goal Models}

\author{Tim Menzies}
\affiliation{\institution{}\city{}\country{}}
\email{timm@ieee.org}

\begin{abstract}
Goal-oriented requirements engineering is rich in notations and
comparatively poor in inference. These notes place a small
artifact---a ${\sim}100$-line stochastic interpreter for
softgoal graphs plus a delta-debugging key extractor---against
that literature. We trace the ``requirements engineering
equation'' from entailment to search, identify what the
artifact should be compared against, list the open problems it
touches and the functionality it adds, discuss how i* models
relate to use cases (and whether ordering subgoals closes that
gap), and assess the value of pairing this inference engine
with LLM-based model extraction. All reference links were
resolved and verified on 2026-08-16.
\end{abstract}

\maketitle

\section{The artifact, in one paragraph}

A goal model is a Python theory: \verb|built <= buy + diy|
declares alternatives, \verb|buy <= vendor * breaks(cheap)|
a conjunction with contribution links; hard goals gate a query
and every softgoal is engaged and labeled in every sampled
world. Worlds are drawn ISAMP-style---one guess per choice
point, restart instead of backtrack, after Crawford and
Baker~\cite{crawford94}. The best of $n_1{=}1000$ worlds (by
distance to the heaven point of maximal quality satisfaction
and minimal leaf purchases) is reduced to a \emph{key}: its
settable labels, minus those every world agrees on, minimized
by Zeller's ddmin~\cite{zeller02} under the test ``30 replays
with this subset stay within $\epsilon$ of the best.'' The
same interpreter replays seeds prudently (believed goals are
settled; disjunctions prefer settled branches). On Horkoff's
seven Kids Help Phone i* models---the corpus of the SHORT
study~\cite{mathew17}---keys of 1--13 labels (1--8\% of model
atoms) steer fresh samples to near the best found world; the
whole corpus runs in about eight seconds.

\section{The requirements engineering equation}

The inference story of RE begins with Jackson's separation of
the world and the machine~\cite{jackson95} and its
crystallization by Zave and Jackson~\cite{zave97}: given
specifications $S$ and domain knowledge $K$, requirements $R$
are met when $S,K \vdash R$. Read literally, the turnstile
makes RE a theorem-proving exercise, and much 1990s-era work
took up exactly that reading. The rejoinder matured inside the
Mylopoulos group: first a revised core ontology of the
requirements problem~\cite{jureta08}, then, in 2010, Techne~\cite{jureta10},
which recast the problem from \emph{proving} $R$ to
\emph{finding} candidate solutions among mandatory and optional
goals under preference---search, not deduction. A parallel
line had long argued the abductive reading: requirements
inference as the generation of assumption sets that explain
desired goals~\cite{menzies96}.

The present artifact is a blunt continuation of that
trajectory: it treats the equation stochastically. Solutions
are not proved or enumerated; they are \emph{sampled}, the best
sample is compressed to the assumptions that matter, and the
compressed answer is validated not by proof but by
\emph{replay}---does re-asserting the key reproduce the
quality? Where Techne asks ``which candidate solutions exist?''
this asks ``which few commitments regenerate a good one on
demand?''

\section{Representation-rich, inference-poor}

The representational lineage is long and healthy: the NFR
framework's softgoals and contribution links~\cite{chung00},
i*'s strategic actors~\cite{yu97}, KAOS's
refinement calculus~\cite{vanlam01}, Tropos and its formal
variant~\cite{fuxman04}, the community-standard cleanup of
iStar~2.0~\cite{dalpiaz16}, lately restated in book form with
a fuller worked example~\cite{dalpiaz24}, and GRL within the
User Requirements Notation~\cite{amyot10}.

Inference over these notations is thinner. The main families:

\begin{itemize}
\item \emph{Qualitative label propagation}, forward and
  backward, from the NFR procedure~\cite{chung00} through its
  formalization~\cite{giorgini02};
\item \emph{SAT-based} minimum-cost satisfiability over goal
  graphs~\cite{sebastiani04}, later generalized to
  constrained goal models solved by Optimization Modulo
  Theories~\cite{nguyen16};
\item \emph{Interactive} evaluation that deliberately keeps a
  human in the propagation loop~\cite{horkoff14};
\item \emph{Metaheuristic search}: SHORT's polynomial-time key
  ranking, evaluated against NSGA-II~\cite{mathew17}.
\end{itemize}

Horkoff and Yu's comparison of satisfaction analysis
techniques~\cite{horkoff12} documents the asymmetry directly:
many notations, few analysis procedures, little tooling, and
results that rarely travel between notations. The extended
systematic mapping of the field~\cite{horkoff19}---246
top-cited GORE papers---quantifies it: proposals and
formalizations dominate, while papers that \emph{evaluate}
analyses remain scarce. That imbalance is an argument for
cheap, executable baselines: an analysis that runs in seconds
on the standard corpus lowers the price of the missing
evaluations.

Our niche in that map is \emph{simulation plus test-driven
minimization}. The minimization side imports a lineage RE has
barely used: ddmin~\cite{zeller02} and its modern probabilistic and
weighted descendants~\cite{probdd24,wdd25}, and---closest in
spirit---%
prediction-preserving input minimization, which delta-debugs an
input to the fragment that keeps an ML model's
answer~\cite{suneja20}. We delta-debug an \emph{assumption set}
to the fragment that keeps a \emph{stochastic simulator's}
answer; to our knowledge that combination is new in RE.

\paragraph{What to compare against.} Four baselines are
natural, in increasing strength: (1)~SHORT itself, same corpus,
same question (done: our keys are 1--8\% of atoms against their
``often just 12\%''); (2)~exact optima from minimum-cost
SAT~\cite{sebastiani04} or OMT~\cite{nguyen16} on the same
models, giving the optimality gap of sampling; (3)~NSGA-II as
the standard many-objective yardstick~\cite{mathew17}; and
(4)~interactive evaluation~\cite{horkoff14} as the human
baseline for ``how many decisions must a stakeholder actually
make?''---which is what a key claims to bound.

\section{Open problems; new functionality}

\emph{Addressed.} (a)~\emph{Trade-space summarization}: rather
than a Pareto front of hundreds of worlds, the output is a
handful of labels a stakeholder can read, audit, and refuse.
(b)~\emph{Validated recommendations}: keys are tested by
replay, so ``do these five things'' arrives with a measured
mean and variance, not a promise. (c)~\emph{Luck versus
choice}: some models' best worlds are partly coin fortune from
help/hurt links; no key can guarantee them, and the ratio of
soft links predicts the gap. Naming that bound---what fraction
of observed quality is \emph{decidable at all}---seems to be
new. (d)~\emph{Scale of machinery}: two files, ${\sim}150$
lines of code plus seven named constants, eight seconds per
corpus; small enough for CI, classrooms, and audit.

\emph{Enabled.} Seeded what-if analysis (assert any partial
belief set, sample the consequences); ablation of stakeholder
positions (drop one commitment from a key, watch the quality
distribution move); prudence checking as a first-class
operation (a claimed optimum that cannot be replayed from its
own labels is an artifact---this check caught a real bug in
our own port).

\section{iStar 2.0 and use cases}

iStar 2.0 and use cases answer different questions. A goal
model is an \emph{intentional variability space}: why actors
want things, which alternatives exist, how qualities trade.
A use case is a \emph{behavioral contract}: actor--system
interaction sequences, main and exceptional flows, pre- and
post-conditions. The bridge has been walked in one direction:
Santander and Castro derive use cases from i*
models~\cite{santander02}, and the User Requirements Notation
pairs GRL goal models with Use Case Maps precisely because
neither view subsumes the other~\cite{amyot10}.

\emph{Does ordering subgoals give use cases?} Partly. Our
kernel is already closer than it looks: a bare body list runs
in order (only the explicit \verb|and| shuffles), so
\emph{sequence is one dialect convention away}, and worlds
would become traces. That yields scenario skeletons: ordered,
branching (via \verb|or|), abortable (a failed demand ends the
trace). What it does not yield is what use cases are for:
actor--system \emph{message} structure, alternative flows
selected by \emph{conditions} rather than dice, interruption
and resumption, and data. Order is necessary but not
sufficient; the URN experience suggests keeping the two views
linked rather than merging them. Recent benchmarks of LLM use
case recovery~\cite{ucrbench25} suggest the behavioral side may
soon be extractable too, which strengthens the two-view
position.

\section{LLMs as model extractors}

Can LLMs read large documents and emit softgoal graphs?
The extraction question is already an active subfield:
comparative studies of LLMs for goal model
extraction~\cite{chen24}, extraction from user
stories~\cite{gunes24}, end-to-end generative goal
modeling~\cite{gengoal25}, sobering evaluations of prompting
limits~\cite{llmgoal26}, and extraction of requirements from
upstream sources such as online user feedback~\cite{feedback25}.
The 2025 systematic review of LLMs in RE~\cite{llm4re25}
(74 studies, 2023--24) finds the field GPT-centric, zero- and
few-shot, evaluated in controlled settings, and---most relevant
here---rarely \emph{integrated into downstream workflows}: the
models extract artifacts that nothing then consumes. The
formal-methods community has begun sketching the missing loop,
a two-way roadmap where LLMs draft formal artifacts and formal
tools check them~\cite{roadmap25,formalsurvey25}. Extraction
\emph{alone} is of limited value---it produces diagrams that
humans must then read, and a 300-node hallucination-prone
diagram is not obviously better than the document it
summarizes.

The added value appears when extraction feeds an inference
engine like this one:

\begin{itemize}
\item \emph{Hallucination linting}: undefined or unlinked atoms
  are mechanical to detect; a graph that cannot be sampled is
  rejected before anyone reads it.
\item \emph{Sanity by simulation}: label histograms over 1000
  worlds expose absurd structure (a ``quality'' that is always
  satisfied, a hard goal that never is) far faster than visual
  inspection.
\item \emph{Audit by key}: the human reviews 5--13 extracted
  decisions---each traceable back to source text---instead of
  the whole graph. The key becomes the interface between the
  LLM's recall and the human's judgment.
\item \emph{Robustness signals}: keys stable across seeds and
  across re-extractions are evidence about the document; keys
  that wobble localize where the LLM guessed.
\end{itemize}

The dialect helps here: the model \emph{is} Python
(\verb|built <= buy + diy|), a friendlier LLM output target
than XML interchange formats, and syntax errors fail fast.
The risk is unchanged from all LLM pipelines---fluent wrong
structure---but a key of ten auditable commitments is a far
smaller attack surface than a diagram, and the replay test
gives every extracted model a falsifiable claim to fail.

\section{Conclusion}

The field's imbalance is real: notations accumulated for three
decades while inference stayed thin. The artifact described
here is one more entry on the thin side---sampling as the
solver, delta debugging as the explainer, replay as the
validator---cheap enough to run anywhere and small enough to
audit. Its natural next experiments are the optimality-gap
comparison against exact solvers, a sequence-extended dialect
evaluated against derived use cases, and an
extract-lint-sample-key pipeline over real documents.

\begin{thebibliography}{36}

\bibitem{jackson95} M.~Jackson. The world and the machine.
\emph{ICSE 1995}. \url{https://doi.org/10.1145/225014.225041}

\bibitem{zave97} P.~Zave, M.~Jackson. Four dark corners of
requirements engineering. \emph{ACM TOSEM} 6(1), 1997.
\url{https://doi.org/10.1145/237432.237434}

\bibitem{jureta08} I.~Jureta, J.~Mylopoulos, S.~Faulkner.
Revisiting the core ontology and problem in requirements
engineering. \emph{RE 2008}.
\url{https://doi.org/10.1109/re.2008.13}

\bibitem{jureta10} I.~Jureta, A.~Borgida, N.~Ernst,
J.~Mylopoulos. Techne: Towards a new generation of requirements
modeling languages. \emph{RE 2010}.
\url{https://doi.org/10.1109/re.2010.24}

\bibitem{menzies96} T.~Menzies. Applications of abduction:
knowledge-level modelling. \emph{Int.\ J.\ Human-Computer
Studies} 45(3), 1996.
\url{https://doi.org/10.1006/ijhc.1996.0054}

\bibitem{chung00} L.~Chung, B.~Nixon, E.~Yu, J.~Mylopoulos.
\emph{Non-Functional Requirements in Software Engineering}.
Springer, 2000.
\url{https://doi.org/10.1007/978-1-4615-5269-7}

\bibitem{yu97} E.~Yu. Towards modelling and reasoning support
for early-phase requirements engineering. \emph{ISRE 1997}.
\url{https://doi.org/10.1109/isre.1997.566873}

\bibitem{vanlam01} A.~van Lamsweerde. Goal-oriented
requirements engineering: a guided tour. \emph{RE 2001}.
\url{https://doi.org/10.1109/isre.2001.948567}

\bibitem{fuxman04} A.~Fuxman, L.~Liu, J.~Mylopoulos,
M.~Pistore, M.~Roveri, P.~Traverso. Specifying and analyzing
early requirements in Tropos. \emph{Requirements Engineering}
9(2), 2004. \url{https://doi.org/10.1007/s00766-004-0191-7}

\bibitem{dalpiaz16} F.~Dalpiaz, X.~Franch, J.~Horkoff. iStar
2.0 language guide, 2016.
\url{https://arxiv.org/abs/1605.07767}

\bibitem{dalpiaz24} F.~Dalpiaz, X.~Franch, J.~Horkoff.
Definition and use of iStar 2.0. In \emph{Social Modeling
Using the i* Framework}, Springer, 2024.
\url{https://doi.org/10.1007/978-3-031-72107-6_5}

\bibitem{amyot10} D.~Amyot, S.~Ghanavati, J.~Horkoff,
G.~Mussbacher, L.~Peyton, E.~Yu. Evaluating goal models within
the goal-oriented requirement language. \emph{Int.\ J.\
Intelligent Systems} 25(8), 2010.
\url{https://doi.org/10.1002/int.20433}

\bibitem{giorgini02} P.~Giorgini, J.~Mylopoulos,
E.~Nicchiarelli, R.~Sebastiani. Reasoning with goal models.
\emph{ER 2002}.
\url{https://doi.org/10.1007/3-540-45816-6_22}

\bibitem{sebastiani04} R.~Sebastiani, P.~Giorgini,
J.~Mylopoulos. Simple and minimum-cost satisfiability for goal
models. \emph{CAiSE 2004}.
\url{https://doi.org/10.1007/978-3-540-25975-6_4}

\bibitem{nguyen16} C.~M.~Nguyen, R.~Sebastiani, P.~Giorgini,
J.~Mylopoulos. Multi-objective reasoning with constrained goal
models. \emph{Requirements Engineering} 23, 2018.
\url{https://doi.org/10.1007/s00766-016-0263-5}

\bibitem{horkoff14} J.~Horkoff, E.~Yu. Interactive goal model
analysis for early requirements engineering.
\emph{Requirements Engineering} 21, 2016.
\url{https://doi.org/10.1007/s00766-014-0209-8}

\bibitem{horkoff12} J.~Horkoff, E.~Yu. Comparison and
evaluation of goal-oriented satisfaction analysis techniques.
\emph{Requirements Engineering} 18, 2013.
\url{https://doi.org/10.1007/s00766-011-0143-y}

\bibitem{mathew17} G.~Mathew, T.~Menzies, N.~Ernst, J.~Klein.
SHORT: shortcutting the search for solutions in goal models,
2017. \url{https://arxiv.org/abs/1702.05568}

\bibitem{crawford94} J.~Crawford, A.~Baker. Experimental
results on the application of satisfiability algorithms to
scheduling problems. \emph{AAAI 1994}.
\url{https://cdn.aaai.org/AAAI/1994/AAAI94-168.pdf}

\bibitem{zeller02} A.~Zeller, R.~Hildebrandt. Simplifying and
isolating failure-inducing input. \emph{IEEE TSE} 28(2), 2002.
\url{https://doi.org/10.1109/32.988498}

\bibitem{probdd24} M.~Zhang et al. Toward a better
understanding of probabilistic delta debugging, 2024.
\url{https://arxiv.org/abs/2408.04735}

\bibitem{suneja20} S.~Suneja, Y.~Zheng, Y.~Zhuang, J.~Laredo,
A.~Morari. Probing model signal-awareness via
prediction-preserving input minimization, 2020.
\url{https://arxiv.org/abs/2011.14934}

\bibitem{santander02} V.~Santander, J.~Castro. Deriving use
cases from organizational modeling. \emph{RE 2002}.
\url{https://doi.org/10.1109/icre.2002.1048503}

\bibitem{ucrbench25} UCRBench: benchmarking LLMs on use case
recovery, 2025. \url{https://arxiv.org/abs/2512.13360}

\bibitem{chen24} A comparative study of large language models
for goal model extraction. \emph{MODELS Companion (SAM) 2024}.
\url{https://doi.org/10.1145/3652620.3686246}

\bibitem{gunes24} Goal model extraction from user stories
using large language models. Springer LNCS, 2024.
\url{https://doi.org/10.1007/978-3-031-70245-7_19}

\bibitem{gengoal25} Generative goal modeling, 2025.
\url{https://arxiv.org/abs/2509.01048}

\bibitem{llmgoal26} Evaluating LLM-based goal extraction in
requirements engineering: prompting strategies and their
limitations, 2026. \url{https://arxiv.org/abs/2604.22207}

\bibitem{horkoff19} J.~Horkoff, F.~Aydemir, E.~Cardoso,
T.~Li, A.~Mat\'e, E.~Paja, M.~Salnitri, L.~Piras,
J.~Mylopoulos, P.~Giorgini. Goal-oriented requirements
engineering: an extended systematic mapping study.
\emph{Requirements Engineering} 24, 2019.
\url{https://doi.org/10.1007/s00766-017-0280-z}

\bibitem{wdd25} R.~Wang et al. WDD: weighted delta debugging.
\emph{ICSE 2025}.
\url{https://doi.org/10.1109/ICSE55347.2025.00071}

\bibitem{llm4re25} Large language models (LLMs) for
requirements engineering (RE): a systematic literature review,
2025. \url{https://arxiv.org/abs/2509.11446}

\bibitem{roadmap25} Formal requirements engineering and large
language models: a two-way roadmap. \emph{Information and
Software Technology} 181, 2025.
\url{https://doi.org/10.1016/j.infsof.2025.107697}

\bibitem{formalsurvey25} A short survey on formalising
software requirements using large language models, 2025.
\url{https://arxiv.org/abs/2506.11874}

\bibitem{feedback25} From online user feedback to
requirements: evaluating large language models, 2025.
\url{https://arxiv.org/abs/2510.23055}

\bibitem{zadeh65} L.~Zadeh. Fuzzy sets. \emph{Information and
Control} 8(3), 1965.
\url{https://doi.org/10.1016/S0019-9958(65)90241-X}

\end{thebibliography}
\end{document}
